% Capítulo: El Dios Todopoderoso (generado desde 09_Omnipotente/09_Omnipotente_capitulo.tex)
% Texto ancla: Génesis 17:1 y 18:14. Clasificación: Identidad; incomunicable (reflejo: fortaleza recibida y dominio de sí).

\section{Introducción: no se alabe el valiente}

El verso que une esta obra nombra tres glorias falsas, y la segunda de ellas es la del poder:

\hebverso{אַל־יִתְהַלֵּל חָכָם בְּחָכְמָתוֹ וְאַל־יִתְהַלֵּל הַגִּבּוֹר בִּגְבוּרָתוֹ\\ הַשְׂכֵּל וְיָדֹעַ אוֹתִי כִּי אֲנִי יְהוָה}
\begin{verso}{Jeremías 9:23–24 [9:22–23]}
Así dijo \SENOR: No se alabe el sabio en su sabiduría, ni en su valentía se alabe el valiente, ni el rico se alabe en sus riquezas. Mas alábese en esto el que se hubiere de alabar: en entenderme y conocerme, que yo soy \SENOR, que hago misericordia, juicio y justicia en la tierra; porque estas cosas quiero, dice \SENOR.
\end{verso}

El ``valiente'' es el \emph{gibbor} (\heb{גִּבּוֹר}), y su ``valentía'', la \emph{gevurá} (\heb{גְּבוּרָה}). Con esa misma raíz la Torá y los profetas llaman a Dios \heb{הָאֵל הַגִּבּוֹר}, ``el Dios poderoso'' (Dt 10:17; Jer 32:18). Jeremías no niega que el poder humano exista; lo que niega es que sea un lugar donde gloriarse. El único \emph{gibbor} digno de gloria es el que hace misericordia, juicio y justicia.

La definición de la que parte este capítulo cabe en una línea: \emph{Dios tiene poder infinito, y nada puede frustrar su voluntad}. Tiene dos mitades. La primera afirma que no hay nada que Dios no pueda hacer; la segunda, que nadie puede impedirle hacer lo que quiere. Job las confiesa juntas al final de su prueba:

\hebverso{יָדַעְתָּ כִּי־כֹל תּוּכָל וְלֹא־יִבָּצֵר מִמְּךָ מְזִמָּה}
\begin{verso}{Job 42:2}
Yo conozco que todo lo puedes, y que no hay pensamiento que se esconda de ti.
\end{verso}

Isaías, por su parte, las convierte en desafío:

\begin{verso}{Isaías 14:27}
Porque \SENOR{} de los ejércitos lo ha determinado, ¿y quién lo impedirá? Y su mano extendida, ¿quién la hará retroceder?
\end{verso}

La palabra ``omnipotencia'' no está en la Biblia. Es el nombre que la teología da a lo que Israel confiesa con un Nombre, \heb{אֵל שַׁדַּי} (\emph{El Shaddai}), y con una pregunta: \heb{הֲיִפָּלֵא מֵיְהוָה דָּבָר}, ``¿hay para \SENOR{} alguna cosa difícil?'' (Gn 18:14). Israel no aprendió este atributo en un tratado; lo aprendió en la tienda de dos ancianos sin hijos y en la orilla de un mar que se abrió.

Esta serie sitúa la omnipotencia entre los atributos de la \emph{identidad}, los que responden a la pregunta ¿quién es Dios? Como todos los atributos, es una perfección eterna de Dios, y es \emph{incomunicable}: solo Dios es \emph{Shaddai}, y ninguna criatura tiene un poder que no haya recibido (Dt 8:17--18). Deja, sin embargo, un \emph{reflejo} en el ser humano, porque Dios ``da esfuerzo al cansado'' (Is 40:29) y su poder ``se perfecciona en la debilidad'' (2~Co 12:9). El ser humano no llega a ser omnipotente; llega a ser fortalecido.

\begin{center}\small\renewcommand{\arraystretch}{1.25}
\begin{tabular}{@{}L{2.6cm}L{5.4cm}L{6.4cm}@{}}
\toprule
\textbf{Mirada} & \textbf{Pregunta} & \textbf{Textos centrales} \\ \midrule
Hebrea & ¿Cómo confiesa Israel el poder de Dios? & Gn 17:1; 18:14; Éx 14--15; Dt 10:17; 32:39; Sal 33; 62; 91; 115; Job 42:2 \\
Judía & ¿Cómo lo leen los sabios de Israel? & \emph{Jaguigá} 12a; Rashi a Gn 17:1; Amidá; \emph{Taanit} 2a; \emph{Yomá} 69b; \emph{Meguilá} 31a; 4Q521 \\
Teológica & ¿Qué afirma la Escritura del poder y qué exige? & Sal 62:11; Jer 32:17; Hab 3:4; Job 26:14; He 1:3; Ro 9:22 \\
Paulina & ¿Qué lugar tiene en Cristo? & Ro 1:16, 20; 4:17--21; 1~Co 1:18--25; 2~Co 6:18; 12:9; Ef 1:19--21; Flp 3:21 \\
\bottomrule
\end{tabular}\end{center}

%===================================================================
\section{Texto ancla: Génesis 17:1 y 18:14}

\subsection{\emph{El Shaddai} a un hombre de noventa y nueve años}

Trece años separan el final de Génesis 16 del comienzo de Génesis 17. Ismael ya es un muchacho, y la promesa de un hijo de Sarai parece haberse disuelto en una solución humana. En ese silencio, Dios habla:

\hebverso{וַיְהִי אַבְרָם בֶּן־תִּשְׁעִים שָׁנָה וְתֵשַׁע שָׁנִים וַיֵּרָא יְהוָה אֶל־אַבְרָם\\ וַיֹּאמֶר אֵלָיו אֲנִי־אֵל שַׁדַּי הִתְהַלֵּךְ לְפָנַי וֶהְיֵה תָמִים}
\begin{verso}{Génesis 17:1–2}
Era Abram de edad de noventa y nueve años, cuando le apareció \SENOR{} y le dijo: Yo soy el Dios Todopoderoso; anda delante de mí y sé perfecto. Y pondré mi pacto entre mí y ti, y te multiplicaré en gran manera.
\end{verso}

\begin{filologia}{Gn 17:1--2}
\heb{אֵל שַׁדַּי} \emph{El Shaddai} & \heb{שׁדי}? & ``Dios Todopoderoso'' (RVR1960). Primera aparición del Nombre. Éx 6:3 lo identifica como el nombre con que Dios se dio a conocer a los patriarcas. La etimología se discute (sección~\ref{sec:shaddai}). \\
\heb{הִתְהַלֵּךְ לְפָנַי} & \heb{הלך} & ``anda delante de mí'': forma reflexiva-iterativa (\emph{hitpael}), ``ve y ven, vive'' bajo la mirada de Dios. \\
\heb{תָמִים} \emph{tamim} & \heb{תמם} & ``entero, íntegro, completo''. La respuesta humana al Todopoderoso no es un acto de fuerza, sino una vida entera. \\
\heb{וְאַרְבֶּה אוֹתְךָ} & \heb{רבה} & ``y te multiplicaré'': el primer contenido del Nombre es fecundidad. \\
\end{filologia}

El Nombre llega en el momento exacto en que el poder humano se ha agotado. Abram tiene noventa y nueve años; Sarai es estéril y ha dejado atrás la edad de concebir (18:11). \emph{El Shaddai} se presenta allí donde nada puede nacer, y lo primero que promete es multiplicar.

\subsection{“¿Hay para Dios alguna cosa difícil?”}

En el capítulo siguiente, tres visitantes llegan a la tienda de Mamre. Uno de ellos anuncia que dentro de un año Sara tendrá un hijo. Sara, que escucha detrás de la puerta, se ríe:

\begin{verso}{Génesis 18:11–13}
Y Abraham y Sara eran viejos, de edad avanzada; y a Sara le había cesado ya la costumbre de las mujeres. Se rió, pues, Sara entre sí, diciendo: ¿Después que he envejecido tendré deleite, siendo también mi señor ya viejo? Entonces \SENOR{} dijo a Abraham: ¿Por qué se ha reído Sara diciendo: ¿Será cierto que he de dar a luz siendo ya vieja?
\end{verso}

La respuesta de Dios toma la forma de una pregunta:

\hebverso{הֲיִפָּלֵא מֵיְהוָה דָּבָר לַמּוֹעֵד אָשׁוּב אֵלֶיךָ כָּעֵת חַיָּה וּלְשָׂרָה בֵן}
\begin{verso}{Génesis 18:14}
¿Hay para Dios alguna cosa difícil? Al tiempo señalado volveré a ti, y según el tiempo de la vida, Sara tendrá un hijo.
\end{verso}

\begin{filologia}{Gn 18:14}
\heb{הֲיִפָּלֵא} \emph{hayippalé} & \heb{פלא} & nifal de \emph{palá}: ``¿será extraordinario, fuera de alcance, demasiado difícil?''. La raíz designa lo que excede la capacidad humana: lo ``maravilloso'' (Sal 139:6) y lo ``difícil'' de juzgar (Dt 17:8). \\
\heb{מֵיְהוָה} \emph{meAdonay} & --- & ``para \SENOR'', literalmente ``desde \SENOR'': medido con la capacidad de Dios. La RVR1960 traduce ``para Dios''. \\
\heb{דָּבָר} \emph{davar} & \heb{דבר} & ``palabra'' y ``cosa''. La pregunta admite las dos lecturas: ¿hay alguna \emph{cosa} difícil?, ¿hay alguna \emph{palabra} suya que no pueda cumplir? \\
\heb{לַמּוֹעֵד} \emph{lammoed} & \heb{יעד} & ``al tiempo señalado'': el poder de Dios también es dueño del tiempo. \\
\heb{כָּעֵת חַיָּה} & \heb{חיה} & ``según el tiempo de la vida'': el plazo de un embarazo, ``cuando el tiempo reviva''. \\
\end{filologia}

La Septuaginta vertió la pregunta como \grc{μὴ ἀδυνατεῖ παρὰ τῷ θεῷ ῥῆμα;}, ``¿será imposible para Dios una palabra?''. Conviene retener la forma griega, porque siglos después el ángel la repetirá a María. La promesa, entre tanto, se cumple al pie de la letra:

\begin{verso}{Génesis 21:1–2}
Visitó \SENOR{} a Sara, como había dicho, e hizo \SENOR{} con Sara como había hablado. Y Sara concibió y dio a Abraham un hijo en su vejez, en el tiempo que Dios le había dicho.
\end{verso}

Tres veces repite el relato que Dios hizo ``como había dicho''. La omnipotencia, en la Torá, no empieza como una idea sobre lo que Dios podría hacer en abstracto; es la fidelidad de una palabra que se cumple en un cuerpo donde la vida ya no era posible.

\subsection{El hilo de \heb{פלא}: lo que es “demasiado” para el hombre}

La raíz de Génesis 18:14 recorre la Biblia hebrea como un hilo. Está en el cántico del mar, donde Dios es \heb{עֹשֵׂה פֶלֶא}, ``hacedor de prodigios'' (Éx 15:11); en el nombre del niño de Isaías, \heb{פֶּלֶא יוֹעֵץ אֵל גִּבּוֹר}, ``Admirable, Consejero, Dios Fuerte'' (Is 9:6 [9:5]); y en la oración de Jeremías, que compra un campo mientras Babilonia sitia Jerusalén:

\hebverso{הִנֵּה אַתָּה עָשִׂיתָ אֶת־הַשָּׁמַיִם וְאֶת־הָאָרֶץ בְּכֹחֲךָ הַגָּדוֹל\\ וּבִזְרֹעֲךָ הַנְּטוּיָה לֹא־יִפָּלֵא מִמְּךָ כָּל־דָּבָר}
\begin{verso}{Jeremías 32:17}
¡Oh Señor \SENOR! he aquí que tú hiciste el cielo y la tierra con tu gran poder, y con tu brazo extendido, ni hay nada que sea difícil para ti.
\end{verso}

Dios le responde con las mismas palabras de la tienda de Mamre:

\begin{verso}{Jeremías 32:27}
He aquí que yo soy \SENOR, Dios de toda carne; ¿habrá algo que sea difícil para mí?
\end{verso}

El hilo pasa, por último, por la confesión de Job, que reconoce haber hablado de \heb{נִפְלָאוֹת מִמֶּנִּי}, ``cosas demasiado maravillosas para mí'' (Job 42:3). Lo que es \emph{pele} para el ser humano ---lo que lo sobrepasa--- no lo es para Dios. En el Nuevo Testamento, el ángel lo anuda con la frase griega de la Septuaginta: \grc{οὐκ ἀδυνατήσει παρὰ τοῦ θεοῦ πᾶν ῥῆμα}.

\begin{verso}{Lucas 1:37}
porque nada hay imposible para Dios.
\end{verso}

\begin{center}\small\renewcommand{\arraystretch}{1.25}
\begin{tabular}{@{}L{2.6cm}L{5.4cm}L{6.4cm}@{}}
\toprule
\textbf{Texto} & \textbf{Forma} & \textbf{Situación} \\ \midrule
Gn 18:14 & \heb{הֲיִפָּלֵא} / LXX \grc{μὴ ἀδυνατεῖ} & una anciana estéril \\
Éx 15:11 & \heb{עֹשֵׂה פֶלֶא} & un pueblo sin salida ante el mar \\
Is 9:6 [9:5] & \heb{פֶּלֶא}, nombre del niño & un reino en tinieblas (9:2 [9:1]) \\
Jer 32:17, 27 & \heb{לֹא־יִפָּלֵא} & una ciudad sitiada \\
Job 42:3 & \heb{נִפְלָאוֹת} & un justo que sufre \\
Lc 1:37 & \grc{οὐκ ἀδυνατήσει} & una virgen \\
\bottomrule
\end{tabular}\end{center}

En cada uno de estos casos, lo ``demasiado difícil'' es precisamente lo que el ser humano no puede producir: un hijo, un camino en el mar, un reino, un futuro, una respuesta al sufrimiento. Allí, y no en abstracto, muestra la Escritura el poder de Dios.

%===================================================================
\section{La palabra}

\subsection{“Todopoderoso” y “Omnipotente” en la RVR1960}\label{sec:shaddai}

La Reina-Valera 1960 traduce el Nombre con dos palabras, ``Todopoderoso'' y ``Omnipotente''. Según el conteo propio con la RVR1960 y la BHS (sin los títulos de sección que agregan los editores), entre ambas suman \textbf{58 veces}. Las 48 del Antiguo Testamento traducen siempre \heb{שַׁדַּי} (\emph{Shaddai}, Strong H7706): 17 como ``Todopoderoso'' y 31 como ``Omnipotente''. Las 10 del Nuevo Testamento traducen \grc{παντοκράτωρ} (\emph{pantokrátōr}, G3841), una vez en Pablo (2~Co 6:18) y nueve en Apocalipsis. La palabra ``omnipotencia'', en cambio, no aparece ninguna vez.

\begin{center}\small\renewcommand{\arraystretch}{1.2}
\begin{tabular}{@{}L{3.2cm}cL{9.6cm}@{}}
\toprule
\textbf{Libro} & \textbf{Veces} & \textbf{Textos} \\ \midrule
Job & 31 & 5:17 a 40:2; nombre preferido en los diálogos \\
Génesis & 6 & 17:1; 28:3; 35:11; 43:14; 48:3; 49:25 \\
Números & 2 & 24:4, 16 (Balaam) \\
Rut & 2 & 1:20, 21 (Noemí) \\
Salmos & 2 & 68:14 [68:15]; 91:1 \\
Ezequiel & 2 & 1:24; 10:5 \\
Éxodo & 1 & 6:3 \\
Isaías & 1 & 13:6 \\
Joel & 1 & 1:15 \\ \midrule
Total AT & 48 & \heb{שַׁדַּי}; siete veces en la forma \heb{אֵל שַׁדַּי} (Gn 17:1; 28:3; 35:11; 43:14; 48:3; Éx 6:3; Ez 10:5) \\
Nuevo Testamento & 10 & \grc{παντοκράτωρ}: 2~Co 6:18; Ap 1:8; 4:8; 11:17; 15:3; 16:7, 14; 19:6, 15; 21:22 \\
\bottomrule
\end{tabular}\end{center}

\precision{conteo propio con la RVR1960 y la BHS etiquetada con Strong; se excluyen los nombres propios que contienen \emph{Shaddai} (Zurisadai y Amisadai, en Números). El etiquetado de la BHS de la biblioteca omite dos casos (Job 24:1; 34:12), que se recuperaron por la RVR1960. Conviene verificarlo en e-Sword.}

La distribución no es neutra. Dos de cada tres veces el Nombre está en Job: es el nombre con que se discute el sufrimiento del justo. Job dice que ``las saetas del Todopoderoso están en mí'' (6:4); Noemí, que ``en grande amargura me ha puesto el Todopoderoso'' (Rt 1:20). El poder de Dios no entra en la Biblia como un dato tranquilizador, sino como la realidad con la que el creyente tiene que habérselas en la prueba.

\subsection{¿Qué significa \emph{Shaddai}?}

La etimología del Nombre sigue sin resolverse. Cuatro propuestas se disputan el campo:

\begin{center}\small\renewcommand{\arraystretch}{1.25}
\begin{tabular}{@{}L{3.0cm}L{6.4cm}L{5.0cm}@{}}
\toprule
\textbf{Propuesta} & \textbf{Base} & \textbf{Textos y fuente} \\ \midrule
``El que basta'' & lectura rabínica: \heb{שֶׁ} (\emph{she-}, ``que'') + \heb{דַּי} (\emph{dai}, ``suficiente''). La Septuaginta vierte \grc{ὁ ἱκανός}, ``el Suficiente'', en Rt 1:20--21 y en Job 21:15; 31:2; 40:2 & \emph{Jaguigá} 12a; \emph{Bereshit Rabá} 46:3; Rashi a Gn 17:1 \\
``El de la montaña'' & acadio \emph{šadû}, ``montaña'': el Dios alto, firme e inconmovible & Cross (1973); Koehler y Baumgartner (2001) \\
``El de los pechos'' & \heb{שַׁד}, ``pecho'': juego de palabras en la bendición de Jacob & Gn 49:25 \\
``El que asola'' & \heb{שׁדד}, ``devastar'': juego de palabras de los profetas & Is 13:6; Jl 1:15 \\
\bottomrule
\end{tabular}\end{center}

Los dos juegos de palabras bíblicos tiran en direcciones opuestas. En Génesis 49:25, Jacob bendice a José ``por el Dios Omnipotente (\heb{שַׁדַּי})\ldots{} con bendiciones de los pechos (\heb{שָׁדַיִם}) y del vientre''; en Isaías 13:6 y Joel 1:15, el día de \SENOR{} ``vendrá como asolamiento (\heb{שֹׁד}) del Todopoderoso (\heb{מִשַּׁדַּי})''. El mismo poder que hace nacer es el que juzga.

\begin{verso}{Génesis 49:25}
Por el Dios de tu padre, el cual te ayudará, por el Dios Omnipotente, el cual te bendecirá con bendiciones de los cielos de arriba, con bendiciones del abismo que está abajo, con bendiciones de los pechos y del vientre.
\end{verso}

\begin{verso}{Isaías 13:6}
Aullad, porque cerca está el día de \SENOR; vendrá como asolamiento del Todopoderoso.
\end{verso}

\precision{los juegos de palabras muestran cómo oyó Israel el Nombre, no de dónde procede la palabra. Fuera de la Biblia, la inscripción de Deir Alla (Jordania, s.~VIII a.\,e.\,c.), que menciona a Balaam hijo de Beor, habla de unos dioses \emph{šdyn} (Hackett, 1980); el dato sugiere que el término era antiguo en la región, pero no fija su sentido.}

\subsection{\emph{El Shaddai} da vida donde no la hay}

En Génesis el Nombre presenta un rasgo constante: aparece \emph{siempre} junto a promesas de descendencia y de vida.

\begin{center}\small\renewcommand{\arraystretch}{1.2}
\begin{tabular}{@{}L{2.4cm}L{7.4cm}L{4.6cm}@{}}
\toprule
\textbf{Texto} & \textbf{Junto al Nombre} & \textbf{Situación} \\ \midrule
Gn 17:1--2 & ``te multiplicaré en gran manera'' & Abram, 99 años; Sarai estéril \\
Gn 28:3 & ``te haga fructificar y te multiplique'' & Jacob huye sin nada \\
Gn 35:11 & ``crece y multiplícate'' & Jacob en Betel \\
Gn 43:14 & ``os dé misericordia delante de aquel varón'' & Benjamín, el último hijo, en peligro \\
Gn 48:3--4 & ``te haré crecer, y te multiplicaré'' & Jacob anciano bendice a José \\
Gn 49:25 & ``bendiciones de los pechos y del vientre'' & bendición de José \\
\bottomrule
\end{tabular}\end{center}

\begin{verso}{Génesis 48:3–4}
y dijo a José: El Dios Omnipotente me apareció en Luz en la tierra de Canaán, y me bendijo, y me dijo: He aquí yo te haré crecer, y te multiplicaré, y te pondré por estirpe de naciones; y daré esta tierra a tu descendencia después de ti por heredad perpetua.
\end{verso}

En la Torá, \emph{El Shaddai} no es primero el Dios que destruye, sino el que \destacar{hace nacer lo que no puede nacer}. Así lo entenderá Pablo cuando diga que Abraham creyó en el Dios ``que da vida a los muertos'' (Ro 4:17; sección~\ref{sec:cristo-omni}).

\subsection{La familia del poder}

El hebreo no dispone de una sola palabra para ``poder''; tiene una familia de términos, y cada uno trae su propia imagen.

\begin{filologia}{La familia hebrea y griega del poder}
\heb{גְּבוּרָה} \emph{gevurá}; \heb{גִּבּוֹר} \emph{gibbor} & \heb{גבר} & fuerza del guerrero, proeza; ``valiente, poderoso''. Dt 10:17; Sal 24:8; Is 9:6 [9:5]; Jer 9:23 [9:22]. \\
\heb{כֹּחַ} \emph{kóaj} & \heb{כחח} & vigor, capacidad de producir. Éx 15:6; Jer 32:17; Sal 147:5; Is 40:26, 29. \\
\heb{עֹז} \emph{oz} & \heb{עזז} & fuerza firme, fortaleza, refugio. Sal 62:11 [62:12]; Hab 3:4. \\
\heb{זְרוֹעַ נְטוּיָה} \emph{zeróa netuyá} & \heb{זרע} & ``brazo extendido'': el poder en acción en la historia. Éx 6:6; Dt 26:8; Sal 136:12; Jer 32:17. \\
\heb{יָד חֲזָקָה} \emph{yad jazaká} & \heb{חזק} & ``mano fuerte''. Éx 13:9; Dt 3:24. \\
\heb{יָכֹל} \emph{yajol} & \heb{יכל} & ``poder, ser capaz''. Job 42:2: \heb{כֹל תּוּכָל}, ``todo lo puedes''. \\
\heb{בָּצַר} \emph{batsar} (nifal) & \heb{בצר} & ``ser impedido, resultar inalcanzable''. Gn 11:6; Job 42:2. \\
\grc{δύναμις}, \grc{δυνατός} & --- & ``poder, capacidad''; ``poderoso para hacer''. Ro 1:16, 20; 4:21; 1~Co 1:18, 24. \\
\grc{παντοκράτωρ} & --- & de \grc{πᾶς}, ``todo'', y \grc{κράτος}, ``dominio'': ``el que lo domina todo''. En la LXX traduce sobre todo \heb{יְהוָה צְבָאוֹת}, ``\SENOR{} de los ejércitos'', y \emph{Shaddai} en Job. \\
\end{filologia}

En la RVR1960 la palabra ``poder(es)'' aparece unas 296 veces (182 en el AT y 114 en el NT, conteo propio), y muchas de ellas se refieren al poder humano. Los términos de la tabla son los que la Escritura aplica de modo característico al poder de Dios, y ninguno es abstracto: hablan de un guerrero, de un brazo, de una mano, de una voz. Israel piensa el poder de Dios como poder \emph{en acción}.

\subsection{Babel y Job: el mismo verbo}

Un verbo poco frecuente une dos textos lejanos entre sí. En Babel, Dios observa a los hombres que levantan la torre:

\hebverso{וְעַתָּה לֹא־יִבָּצֵר מֵהֶם כֹּל אֲשֶׁר יָזְמוּ לַעֲשׂוֹת}
\begin{verso}{Génesis 11:6}
Y dijo \SENOR: He aquí el pueblo es uno, y todos éstos tienen un solo lenguaje; y han comenzado la obra, y nada les hará desistir ahora de lo que han pensado hacer.
\end{verso}

Job, al final de su prueba, recurre al mismo verbo y a la misma raíz del ``pensar'':

\hebverso{יָדַעְתָּ כִּי־כֹל תּוּכָל וְלֹא־יִבָּצֵר מִמְּךָ מְזִמָּה}

\begin{filologia}{Gn 11:6 y Job 42:2}
\heb{יִבָּצֵר} \emph{yibbatser} & \heb{בצר} & ``ser cortado, quedar fuera de alcance''. En Babel, a los hombres ``nada se les hará inalcanzable''; en Job, a Dios ``ningún propósito se le hace inalcanzable''. \\
\heb{יָזְמוּ} / \heb{מְזִמָּה} & \heb{זמם} & ``planear'' / ``plan, propósito'' (RVR1960: ``lo que han pensado'', ``pensamiento''). La LXX de Job 42:2 traduce el sentido: \grc{ἀδυνατεῖ δέ σοι οὐθέν}, ``nada te es imposible''. \\
\end{filologia}

La humanidad que pretende ser todopoderosa acaba dispersa (Gn 11:8); el hombre que se rinde ante el Todopoderoso acaba viéndolo: ``De oídas te había oído; mas ahora mis ojos te ven'' (Job 42:5). La omnipotencia no es un atributo que el ser humano pueda disputarle a Dios sin destruirse.

%===================================================================
\section{El Nombre que revela la omnipotencia}

En la Escritura el nombre de Dios no es una etiqueta, sino la revelación de lo que él es (Éx 33:18--19; 34:5--7). La omnipotencia está ligada a tres nombres: \emph{El Shaddai}, el nombre de los patriarcas; \emph{\SENOR{} de los ejércitos}, el nombre de los profetas; y \emph{El Gibbor}, el ``Dios Fuerte'' de Isaías. El Nuevo Testamento los recoge en un solo título griego, \grc{παντοκράτωρ}.

\subsection{“Yo soy el Dios Todopoderoso”}

\emph{El Shaddai} es el nombre con que Dios se presentó a Abraham, a Isaac y a Jacob:

\hebverso{וָאֵרָא אֶל־אַבְרָהָם אֶל־יִצְחָק וְאֶל־יַעֲקֹב בְּאֵל שַׁדָּי וּשְׁמִי יְהוָה לֹא נוֹדַעְתִּי לָהֶם}
\begin{verso}{Éxodo 6:3}
Y aparecí a Abraham, a Isaac y a Jacob como Dios Omnipotente, mas en mi nombre \SENOR{} no me di a conocer a ellos.
\end{verso}

\begin{filologia}{Éx 6:3}
\heb{בְּאֵל שַׁדָּי} \emph{beEl Shadday} & \heb{שׁדי}? & ``como Dios Omnipotente'': la \emph{bet} de esencia, ``en calidad de''. Es la forma pausal de \heb{שַׁדַּי}. \\
\heb{וּשְׁמִי יְהוָה} & \heb{שׁם} & ``y mi nombre \SENOR''. Los patriarcas usaron el tetragrama (Gn 12:8; 15:7), pero no conocieron por experiencia lo que el Nombre revela: la redención del éxodo (Éx 6:6--7). \\
\heb{נוֹדַעְתִּי} & \heb{ידע} & ``me di a conocer'' (nifal de \emph{yadá}): conocer por experiencia, no solo de oídas. \\
\end{filologia}

El versículo distingue dos etapas. \emph{El Shaddai} es el Dios que promete y hace nacer (sección~\ref{sec:shaddai}); \SENOR{} es el Dios que cumple y redime con ``brazo extendido'' (Éx 6:6). La omnipotencia se reveló primero como promesa, en la tienda de los patriarcas, y después como liberación, en Egipto.

\subsection{“\SENOR{} de los ejércitos”}

El nombre más frecuente del poder de Dios en los profetas es \heb{יְהוָה צְבָאוֹת} (\emph{Adonay Tsevaot}), ``\SENOR{} de los ejércitos''. Según el conteo propio con la RVR1960 y la BHS, aparece unas \textbf{285 veces} en el Antiguo Testamento: 82 en Jeremías, 62 en Isaías, 53 en Zacarías, 24 en Malaquías, 14 en Hageo, 15 en los Salmos y el resto en Samuel, Reyes, Crónicas y otros profetas. En la Torá no figura ni una vez; su primera mención está en Silo, adonde Elcana subía ``para adorar y para ofrecer sacrificios a \SENOR{} de los ejércitos'' (1~S 1:3).

El nombre reaparece en el enfrentamiento de David con Goliat, donde el poder humano y el poder de Dios se miden cara a cara:

\hebverso{אַתָּה בָּא אֵלַי בְּחֶרֶב וּבַחֲנִית וּבְכִידוֹן\\ וְאָנֹכִי בָא־אֵלֶיךָ בְּשֵׁם יְהוָה צְבָאוֹת אֱלֹהֵי מַעַרְכוֹת יִשְׂרָאֵל}
\begin{verso}{1 Samuel 17:45}
Entonces dijo David al filisteo: Tú vienes a mí con espada y lanza y jabalina; mas yo vengo a ti en el nombre de \SENOR{} de los ejércitos, el Dios de los escuadrones de Israel, a quien tú has provocado.
\end{verso}

Y es también el nombre que los serafines proclaman en la visión de Isaías:

\begin{verso}{Isaías 6:3}
Y el uno al otro daba voces, diciendo: Santo, santo, santo, \SENOR{} de los ejércitos; toda la tierra está llena de su gloria.
\end{verso}

\begin{filologia}{1~S 17:45 e Is 6:3}
\heb{צְבָאוֹת} \emph{tsevaot} & \heb{צבא} & ``ejércitos, huestes''. El término abarca los ejércitos de Israel (\heb{מַעַרְכוֹת יִשְׂרָאֵל}, 1~S 17:45), las huestes celestiales (1~R 22:19) y los astros (Is 40:26: ``él saca y cuenta su ejército''). Dios es el Señor de toda fuerza, visible e invisible. \\
\heb{בְּשֵׁם} \emph{beshem} & \heb{שׁם} & ``en el nombre'': David no opone su fuerza a la de Goliat, sino el Nombre. \\
LXX & --- & la Septuaginta traduce el título de tres maneras: lo translitera, \grc{κύριος σαβαωθ} (Is 6:3; 1~S 17:45); lo vierte \grc{κύριος παντοκράτωρ}, ``Señor Todopoderoso'' (2~S 7:8); o \grc{κύριος τῶν δυνάμεων}, ``Señor de los poderes'' (Sal 24:10 [23:10 LXX]). \\
\end{filologia}

Este nombre une el poder con la presencia. El salmo 46 lo repite como estribillo: ``\SENOR{} de los ejércitos está con nosotros; nuestro refugio es el Dios de Jacob'' (Sal 46:7, 11 [46:8, 12]). Zacarías, a su vez, lo emplea para precisar de qué clase es ese poder: ``No con ejército, ni con fuerza, sino con mi Espíritu, ha dicho \SENOR{} de los ejércitos'' (Zac 4:6). El Señor de los ejércitos no depende de ejércitos.

\subsection{“El Dios Fuerte”}

El tercer nombre es \heb{אֵל גִּבּוֹר} (\emph{El Gibbor}), ``Dios Fuerte''. Isaías lo da al niño que nacerá (Is 9:6 [9:5]) y, un capítulo después, lo aplica a Dios mismo:

\hebverso{שְׁאָר יָשׁוּב שְׁאָר יַעֲקֹב אֶל־אֵל גִּבּוֹר}
\begin{verso}{Isaías 10:21}
El remanente volverá, el remanente de Jacob volverá al Dios fuerte.
\end{verso}

El título que en 10:21 designa al Dios al que vuelve el remanente es, en 9:6 [9:5], el nombre del hijo de David. La Torá ya había llamado a Dios \heb{הָאֵל הַגִּבּוֹר} (Dt 10:17), y Jeremías repite la fórmula: ``Dios grande, poderoso, \SENOR{} de los ejércitos es su nombre'' (Jer 32:18). En estos textos los tres nombres de la omnipotencia llegan a tocarse.

\subsection{El Nombre del poder en Cristo, leído con Pablo}

Pablo usa dos veces los nombres del poder, y en ambos casos los toma de la Septuaginta. En Romanos 9 cita Isaías 1:9 con el título transliterado:

\begin{verso}{Romanos 9:29}
Y como antes dijo Isaías: Si el Señor de los ejércitos no nos hubiera dejado descendencia, como Sodoma habríamos venido a ser, y a Gomorra seríamos semejantes.
\end{verso}

En 2~Corintios 6:18 cita la promesa de Natán con el título traducido: ``dice el Señor Todopoderoso'' (\grc{κύριος παντοκράτωρ}, de 2~S 7:8). El primer texto muestra el poder que preserva un resto, una \grc{σπέρμα}, ``descendencia'', como la que \emph{El Shaddai} prometió a Abraham; el segundo, el poder que adopta, porque el Todopoderoso se hace Padre de ``hijos e hijas''. Y el Apocalipsis, al cantar la visión de Isaías 6, pone en boca de los seres vivientes el título griego: ``Santo, santo, santo es el Señor Dios Todopoderoso'' (Ap 4:8). El \emph{Tsevaot} de los serafines es el \emph{Pantokrátōr} del trono del Cordero.

\begin{center}\small\renewcommand{\arraystretch}{1.25}
\begin{tabular}{@{}L{3.4cm}L{3.6cm}L{3.0cm}L{4.4cm}@{}}
\toprule
\textbf{Nombre o título} & \textbf{Hebreo / griego} & \textbf{Texto} & \textbf{Qué revela} \\ \midrule
Dios Todopoderoso, Dios Omnipotente & \heb{אֵל שַׁדַּי}; LXX \grc{ὁ ἱκανός} (Rt) & Gn 17:1; Éx 6:3 & el Dios que basta y hace nacer lo imposible \\
\SENOR{} de los ejércitos & \heb{יְהוָה צְבָאוֹת}; \grc{κύριος σαβαωθ} & 1~S 17:45; Is 6:3; Ro 9:29 & el Señor de toda fuerza, presente con su pueblo \\
Dios Fuerte & \heb{אֵל גִּבּוֹר} & Is 9:6 [9:5]; 10:21 & el Dios al que vuelve el resto; nombre del hijo de David \\
Señor Todopoderoso & \grc{κύριος παντοκράτωρ} & 2~Co 6:18; Ap 4:8 & el que lo domina todo se hace Padre \\
\bottomrule
\end{tabular}\end{center}

\precision{el midrash \emph{Shemot Rabá} 3:6 pone en boca de Dios un principio sobre sus nombres: ``según mis obras soy llamado''. Entre otras cosas, explica que Dios se llama \emph{Tsevaot} cuando hace guerra contra los impíos, y \emph{El Shadday} cuando suspende el castigo por los pecados del hombre. Es una lectura rabínica: une el nombre de los ejércitos con el poder que juzga, y el de \emph{Shaddai} con el poder que se contiene, como en \emph{Yomá} 69b (sección~\ref{sec:judaismo-omni}).}

\precision{conteo propio de ``\SENOR{} de los ejércitos'' (y ``Dios de los ejércitos'') con la RVR1960 y la BHS; los números redondos que dan las concordancias (alrededor de 285) dependen de si se cuentan las fórmulas con \heb{אֱלֹהִים} y \heb{אֲדֹנָי}. En el Nuevo Testamento el título transliterado aparece solo en Ro 9:29 y Stg 5:4 (\grc{κύριος σαβαώθ}). Conviene verificarlo en e-Sword.}

%===================================================================
\section{Definición}

Reunidos los textos, la omnipotencia puede definirse así:

\begin{center}
\begin{tcolorbox}[colback=papiro,colframe=azul,boxrule=1pt,arc=3pt,width=0.92\linewidth]
La omnipotencia de Dios es la perfección eterna por la cual Él \destacar{puede hacer todo lo que quiere}, sin esfuerzo y sin desgaste, de modo que \destacar{nada puede frustrar su voluntad}; un poder que obra siempre conforme a su naturaleza santa, sabia y fiel (Gn 18:14; Job 42:2; Sal 115:3; Is 40:28; 46:10; 2~Ti 2:13).
\end{tcolorbox}
\end{center}

Tres rasgos sostienen esta definición. El poder de Dios es, ante todo, \emph{sin límite}: nada es ``demasiado'' para él, ni crear, ni sostener, ni resucitar (Gn 18:14; Jer 32:17). Es también \emph{sin oposición}, porque ``no hay quien detenga su mano, y le diga: ¿Qué haces?'' (Dn 4:35 [4:32]; Is 14:27). Y es, por último, \emph{sin desgaste}: Dios ``no desfallece, ni se fatiga con cansancio'' (Is 40:28), y precisamente por eso puede dar fuerzas al que no las tiene.

\subsection{“Todo lo que quiso ha hecho”}

El salmo 115 contrapone el Dios de Israel a los ídolos de las naciones:

\hebverso{וֵאלֹהֵינוּ בַשָּׁמָיִם כֹּל אֲשֶׁר־חָפֵץ עָשָׂה}
\begin{verso}{Salmo 115:3}
Nuestro Dios está en los cielos; todo lo que quiso ha hecho.
\end{verso}

De los ídolos se dice que ``manos tienen, mas no palpan; tienen pies, mas no andan'' (v.~7): poseen la forma del poder sin el poder. El Dios de Israel, en cambio, no tiene forma visible, y sin embargo ``todo lo que quiso ha hecho''. El verbo \heb{חָפֵץ} (\emph{jafets}) significa ``querer, complacerse'', y con él el salmista ata el poder de Dios a su \emph{querer}: no se trata de una fuerza ciega, sino de una voluntad que se cumple. Isaías dice lo mismo desde el otro lado, el del plan:

\begin{verso}{Isaías 46:10}
que anuncio lo por venir desde el principio, y desde la antigüedad lo que aún no era hecho; que digo: Mi consejo permanecerá, y haré todo lo que quiero.
\end{verso}

\heb{עֲצָתִי תָקוּם וְכָל־חֶפְצִי אֶעֱשֶׂה}: ``mi consejo se mantendrá en pie, y todo mi deseo haré''. La omnipotencia garantiza que el consejo de Dios no se quede en proyecto.

\subsection{Lo que Dios “no puede”}

La misma Escritura afirma, sin embargo, que hay cosas que Dios no puede hacer:

\begin{verso}{2 Timoteo 2:13}
Si fuéremos infieles, él permanece fiel; él no puede negarse a sí mismo.
\end{verso}

\begin{verso}{Tito 1:2}
en la esperanza de la vida eterna, la cual Dios, que no miente, prometió desde antes del principio de los siglos.
\end{verso}

Que Dios no pueda mentir ni negarse a sí mismo no limita su poder; lo define. Un poder capaz de traicionarse a sí mismo sería debilidad, no fuerza. La tradición judía llegó a la misma idea por la vía de la filosofía: Maimónides enseña que lo imposible tiene una naturaleza estable que ningún agente produce, y que la incapacidad de hacer lo contradictorio ---por ejemplo, un cuadrado cuya diagonal sea igual a su lado--- no supone falta de poder en Dios (\emph{Guía de perplejos} III, 15). La omnipotencia, por tanto, no es la capacidad de hacer absurdos, sino la de hacer todo lo que Dios, siendo quien es, quiere hacer.

\precision{2~Timoteo y Tito pertenecen a las cartas pastorales, cuya autoría paulina se discute. La misma idea se halla en cartas indiscutidas (Ro 3:3--4: ``¿Su incredulidad habrá hecho nula la fidelidad de Dios? De ninguna manera'') y en la Torá: ``Dios no es hombre, para que mienta'' (Nm 23:19).}

%===================================================================
\section{La omnipotencia en la Torá}

\subsection{En el principio: el poder de la palabra}

El primer acto de poder que narra la Biblia es una palabra:

\hebverso{וַיֹּאמֶר אֱלֹהִים יְהִי אוֹר וַיְהִי־אוֹר}
\begin{verso}{Génesis 1:3}
Y dijo Dios: Sea la luz; y fue la luz.
\end{verso}

Diez veces repite Génesis 1 ``y dijo Dios''. No hay una materia previa que Dios deba vencer, ni otros dioses contra los que luchar, ni esfuerzo alguno. El verbo \heb{בָּרָא} (\emph{bará}, ``crear'', Gn 1:1) solo tiene a Dios por sujeto en toda la Biblia hebrea, y entre el \heb{יְהִי} (``sea'') y el \heb{וַיְהִי} (``y fue'') no media distancia: la palabra de Dios y la realidad coinciden. De ahí que la pregunta de Génesis 18:14 pueda leerse como ``¿hay alguna \emph{palabra} demasiado difícil para \SENOR?''. Quien creó con la palabra no puede ser vencido por ninguna de sus palabras.

\subsection{El Dios grande, poderoso y temible}

En su último discurso, Moisés resume a Dios con tres adjetivos que los judíos siguen repitiendo hoy en la oración:

\hebverso{הָאֵל הַגָּדֹל הַגִּבֹּר וְהַנּוֹרָא אֲשֶׁר לֹא־יִשָּׂא פָנִים וְלֹא יִקַּח שֹׁחַד}
\begin{verso}{Deuteronomio 10:17}
Porque \SENOR{} vuestro Dios es Dios de dioses y Señor de señores, Dios grande, poderoso y temible, que no hace acepción de personas, ni toma cohecho.
\end{verso}

El versículo siguiente precisa en qué emplea Dios ese poder: ``hace justicia al huérfano y a la viuda; que ama también al extranjero dándole pan y vestido'' (10:18). El \emph{gibbor} de la Torá pone su fuerza al servicio de los débiles, y los sabios harán de este enlace una regla (sección~\ref{sec:judaismo-omni}).

En el cántico final de Moisés, el poder de Dios alcanza el límite mismo de la vida y la muerte:

\begin{verso}{Deuteronomio 32:39}
Ved ahora que yo, yo soy, y no hay dioses conmigo; yo hago morir, y yo hago vivir; yo hiero, y yo sano; y no hay quien pueda librar de mi mano.
\end{verso}

\heb{אֲנִי אָמִית וַאֲחַיֶּה}, ``yo hago morir y hago vivir'': el poder sobre la muerte pertenece solo a Dios, y la Torá lo enuncia en presente. La tradición judía vio aquí una alusión a la resurrección, el mismo poder que la Amidá bendice cada día (sección~\ref{sec:judaismo-omni}).

\subsection{“¿Acaso se ha acortado la mano?”}

En el desierto, Moisés duda de que Dios pueda dar carne a seiscientos mil hombres, y recibe por respuesta otra pregunta:

\begin{verso}{Números 11:23}
Entonces \SENOR{} respondió a Moisés: ¿Acaso se ha acortado la mano de \SENOR? Ahora verás si se cumple mi palabra, o no.
\end{verso}

\heb{הֲיַד יְהוָה תִּקְצָר}: ``¿se acortará la mano de \SENOR?''. La mano corta, que no llega a lo que quiere alcanzar, es la imagen hebrea de la impotencia, y los profetas la retoman para negar que Dios sea incapaz de salvar (Is 50:2; 59:1). Una vez más, el poder se mide por el cumplimiento de la palabra: \heb{הֲיִקְרְךָ דְבָרִי}, ``si te sucederá mi palabra''.

\subsection{“Él te da el poder”}

La Torá advierte, además, contra el uso humano del poder:

\begin{verso}{Deuteronomio 8:17–18}
y digas en tu corazón: Mi poder y la fuerza de mi mano me han traído esta riqueza. Sino acuérdate de \SENOR{} tu Dios, porque él te da el poder para hacer las riquezas, a fin de confirmar su pacto que juró a tus padres, como en este día.
\end{verso}

\heb{כֹּחִי וְעֹצֶם יָדִי}, ``mi poder y la fuerza de mi mano'', es la frase del orgullo. Frente a ella, la Torá responde que todo \heb{כֹּחַ} humano es recibido. Es la enseñanza que retomará Jeremías 9:23: el valiente no puede gloriarse en una valentía que no se dio a sí mismo.

%===================================================================
\section{La omnipotencia en los Salmos}

\subsection{“De Dios es el poder”}

El salmo 62 es una meditación sobre la confianza, y al final el salmista condensa lo que ha aprendido:

\hebverso{אַחַת דִּבֶּר אֱלֹהִים שְׁתַּיִם־זוּ שָׁמָעְתִּי כִּי עֹז לֵאלֹהִים\\ וּלְךָ־אֲדֹנָי חָסֶד כִּי־אַתָּה תְשַׁלֵּם לְאִישׁ כְּמַעֲשֵׂהוּ}
\begin{verso}{Salmo 62:11–12 [62:12–13]}
Una vez habló Dios; dos veces he oído esto: que de Dios es el poder, y tuya, oh Señor, es la misericordia; porque tú pagas a cada uno conforme a su obra.
\end{verso}

\begin{filologia}{Sal 62:11--12 (62:12--13)}
\heb{אַחַת}\ldots{} \heb{שְׁתַּיִם} & --- & ``una\ldots{} dos'': paralelismo numérico, recurso de la sabiduría hebrea (cf. Pr 6:16; Am 1:3). Una sola palabra de Dios contiene dos verdades. \\
\heb{עֹז} \emph{oz} & \heb{עזז} & ``poder, fortaleza''. \\
\heb{חֶסֶד} \emph{jésed} & \heb{חסד} & ``misericordia'', amor leal del pacto (ver el capítulo sobre la misericordia). \\
\end{filologia}

El paralelismo junta lo que el lector tendería a separar. \emph{Oz} y \emph{jésed} son las dos cosas que el salmista oyó en una sola palabra de Dios. Un poder sin misericordia aterra, y una misericordia sin poder no salva. Dios habla una vez, y el trueno de su poder se oye en mil colinas (Pink, 2020, cap. 9).

\subsection{“Él dijo, y fue hecho”}

El salmo 33 vuelve a cantar, en verso, el relato de Génesis 1:

\hebverso{בִּדְבַר יְהוָה שָׁמַיִם נַעֲשׂוּ וּבְרוּחַ פִּיו כָּל־צְבָאָם\\ כִּי הוּא אָמַר וַיֶּהִי הוּא־צִוָּה וַיַּעֲמֹד}
\begin{verso}{Salmo 33:6–9}
Por la palabra de \SENOR{} fueron hechos los cielos, y todo el ejército de ellos por el aliento de su boca. Él junta como montón las aguas del mar; él pone en depósitos los abismos. Tema a \SENOR{} toda la tierra; teman delante de él todos los habitantes del mundo. Porque él dijo, y fue hecho; él mandó, y existió.
\end{verso}

\heb{הוּא אָמַר וַיֶּהִי}, ``él dijo y fue'', reproduce la estructura de Génesis 1. La consecuencia que extrae el salmista es el temor, porque el mismo poder que juntó las aguas del mar está frente a ``todos los habitantes del mundo''. El salmo 147, en cambio, añade la ternura:

\begin{verso}{Salmo 147:4–5}
Él cuenta el número de las estrellas; a todas ellas llama por sus nombres. Grande es el Señor nuestro, y de mucho poder; y su entendimiento es infinito.
\end{verso}

Un verso antes, el salmo dice que este mismo Dios ``sana a los quebrantados de corazón, y venda sus heridas'' (v.~3). Quien nombra las estrellas una por una conoce también las heridas de cada persona. Isaías recurre a la misma imagen para consolar a un pueblo cansado: ``él saca y cuenta su ejército; a todas llama por sus nombres; ninguna faltará; tal es la grandeza de su fuerza, y el poder de su dominio'' (Is 40:26).

\subsection{La voz, el brazo, el guerrero y la sombra}

Los salmos expresan el poder de Dios con imágenes concretas:

\begin{center}\small\renewcommand{\arraystretch}{1.25}
\begin{tabular}{@{}L{1.9cm}L{5.6cm}L{6.9cm}@{}}
\toprule
\textbf{Imagen} & \textbf{Texto} & \textbf{Qué dice del poder} \\ \midrule
La voz & Sal 29:4: \heb{קוֹל־יְהוָה בַּכֹּחַ}, ``voz de \SENOR{} con potencia'' & siete veces ``voz de \SENOR'' sobre las aguas, los cedros y el desierto: el trueno es su palabra \\
El brazo & Sal 89:13 [89:14]: \heb{לְךָ זְרוֹעַ עִם־גְּבוּרָה}, ``tuyo es el brazo potente'' & el poder que creó el norte y el sur (v.~12) sostiene la promesa a David (vv. 28--37) \\
El guerrero & Sal 24:8: \heb{יְהוָה עִזּוּז וְגִבּוֹר}, ``\SENOR{} el fuerte y valiente'' & el Rey de gloria entra por las puertas: el poder es realeza \\
La sombra & Sal 91:1: \heb{בְּצֵל שַׁדַּי יִתְלוֹנָן}, ``bajo la sombra del Omnipotente'' & el Nombre más alto se vuelve refugio \\
\bottomrule
\end{tabular}\end{center}

De todas ellas, la última es la más inesperada. El salmo 91 toma el Nombre de Génesis 17 y lo convierte en lugar de descanso:

\hebverso{יֹשֵׁב בְּסֵתֶר עֶלְיוֹן בְּצֵל שַׁדַּי יִתְלוֹנָן}
\begin{verso}{Salmo 91:1–2}
El que habita al abrigo del Altísimo morará bajo la sombra del Omnipotente. Diré yo a \SENOR: Esperanza mía, y castillo mío; mi Dios, en quien confiaré.
\end{verso}

\heb{יִתְלוֹנָן} (\emph{yitlonán}, de \heb{לון}) significa ``pasar la noche'': el creyente duerme a la sombra de \emph{Shaddai}. El poder que podría aplastar es el que cubre.

%===================================================================
\section{Teología: el poder que da vida a todos los atributos}

\subsection{El poder que hace eficaces las demás perfecciones}
Los salmos han mostrado que la Escritura nunca presenta aislado el poder de Dios: el salmista oyó que ``de Dios es el poder'', y en esa misma palabra oyó que ``tuya, oh Señor, es la misericordia'' (Sal 62:11--12). La teología ha extraído de ahí una consecuencia: la omnipotencia no es un atributo más en la lista, sino el que vuelve eficaces a los demás. Stephen Charnock lo formuló con una comparación que Pink recoge al abrir su capítulo: así como la santidad es la belleza de todos los atributos de Dios, el poder es lo que da \emph{vida y acción} a todas sus perfecciones. Sin poder, la misericordia sería ``una débil lástima''; las promesas, ``un sonido vacío''; y las amenazas, ``un mero espantapájaros'' (Pink, 2020, cap. 9). Se entiende así que Jeremías, cuando debe creer una promesa que parece imposible, empiece por el poder: ``¡Oh Señor \SENOR! he aquí que tú hiciste el cielo y la tierra con tu gran poder, y con tu brazo extendido, ni hay nada que sea difícil para ti'' (Jer 32:17). La pregunta de Génesis 18:14 se ha convertido en confesión.

\subsection{Un poder que le pertenece por naturaleza}
El poder de los hombres es prestado: ``acuérdate de \SENOR{} tu Dios, porque él te da el poder para hacer las riquezas'' (Dt 8:18). El de Dios, en cambio, no viene de nadie ni descansa en nada fuera de él. Spurgeon lo expresó con una imagen que Pink cita: Dios ``se sienta en un trono sin columnas que lo sostengan'' (Pink, 2020, cap. 9). Por eso la Escritura puede emplear ``el Poder'' como nombre de Dios, y Jesús habla del Hijo del Hombre ``sentado a la diestra del poder'' (Mr 14:62). La observación de Charnock de que Dios y su poder son tan inseparables que pueden intercambiarse coincide con un uso judío que se examinará en la sección sobre el judaísmo (sección~\ref{sec:judaismo-omni}).

\subsection{El poder escondido}
Lo que Dios muestra de su poder es siempre menos de lo que posee. Habacuc lo afirma en medio de una teofanía que hace temblar los montes:

\hebverso{וְנֹגַהּ כָּאוֹר תִּהְיֶה קַרְנַיִם מִיָּדוֹ לוֹ וְשָׁם חֶבְיוֹן עֻזֹּה}
\begin{verso}{Habacuc 3:4}
Y el resplandor fue como la luz; rayos brillantes salían de su mano, y allí estaba escondido su poder.
\end{verso}

\heb{חֶבְיוֹן} (\emph{jevyón}, de \heb{חבה}, ``esconder'') es el ``escondite''. El resplandor brota de su mano y, sin embargo, el profeta dice que allí el poder estaba \emph{oculto}: aun la teofanía es un velo. Como observa Pink, lo que vemos del poder de Dios oculta más de lo que revela (2020, cap. 9). Job había llegado a la misma conclusión tras recorrer la creación entera: ``He aquí, estas cosas son sólo los bordes de sus caminos; ¡y cuán leve es el susurro que hemos oído de él! Pero el trueno de su poder, ¿quién lo puede comprender?'' (Job 26:14). La omnipotencia no se mide por lo que hemos visto, porque lo que hemos visto es apenas el borde.

\subsection{Creación, preservación, gobierno y juicio}
Ese poder escondido se deja ver, con todo, en cuatro escenarios. Se ve en la \emph{creación}, hecha de la nada y por la palabra: ``él dijo, y fue hecho'' (Sal 33:9). Se ve en la \emph{preservación}, porque el mundo no subsiste por sí mismo: Dios pone límite al mar, ``Hasta aquí llegarás, y no pasarás adelante'' (Job 38:11), y el Hijo ``sustenta todas las cosas con la palabra de su poder'' (He 1:3). Se ve en el \emph{gobierno}, pues aun Satanás actúa solo dentro del límite que Dios le fija (Job 1--2), y las aguas que alzan su voz quedan por debajo del \SENOR{} que ``en las alturas es más poderoso'' (Sal 93:3--4). Y se ve, finalmente, en el \emph{juicio}, manifestado en el diluvio, en Sodoma y en el mar Rojo, donde Dios soportó ``con mucha paciencia los vasos de ira'' para mostrar su poder (Ro 9:22).

\subsection{Temblar, adorar y confiar}
De este poder se siguen tres respuestas, que Pink señala al final de su capítulo (2020, cap. 9). La primera es temblar, y el salmo 2 advierte a los reyes de la tierra que sirvan a \SENOR{} con temor (Sal 2:11--12). La segunda es adorar, como Israel a la orilla del mar: ``¿Quién como tú, oh \SENOR, entre los dioses? ¿Quién como tú, magnífico en santidad, terrible en maravillosas hazañas, hacedor de prodigios?'' (Éx 15:11). La tercera es confiar, porque ninguna oración resulta demasiado difícil de responder para el que todo lo puede: ``\SENOR{} es mi luz y mi salvación; ¿de quién temeré? \SENOR{} es la fortaleza de mi vida; ¿de quién he de atemorizarme?'' (Sal 27:1). Las tres desembocan en la doxología de Efesios 3:20--21, que se comentará en la sección paulina.

%===================================================================
\section{Un encuentro en la historia de Israel: el mar Rojo}

\subsection{Entre el mar y los carros}

Israel aprendió la omnipotencia como pueblo en el paso del mar. Recién salido de Egipto, acampa frente a las aguas mientras a sus espaldas avanzan los carros de Faraón. No tiene armas, ni salida, ni plan. Entonces habla Moisés:

\hebverso{אַל־תִּירָאוּ הִתְיַצְבוּ וּרְאוּ אֶת־יְשׁוּעַת יְהוָה\\ יְהוָה יִלָּחֵם לָכֶם וְאַתֶּם תַּחֲרִישׁוּן}
\begin{verso}{Éxodo 14:13–14}
Y Moisés dijo al pueblo: No temáis; estad firmes, y ved la salvación que \SENOR{} hará hoy con vosotros; porque los egipcios que hoy habéis visto, nunca más para siempre los veréis. \SENOR{} peleará por vosotros, y vosotros estaréis tranquilos.
\end{verso}

\begin{verso}{Éxodo 14:21}
Y extendió Moisés su mano sobre el mar, e hizo \SENOR{} que el mar se retirase por recio viento oriental toda aquella noche; y volvió el mar en seco, y las aguas quedaron divididas.
\end{verso}

\begin{verso}{Éxodo 14:31}
Y vio Israel aquel grande hecho que \SENOR{} ejecutó contra los egipcios; y el pueblo temió a \SENOR, y creyeron a \SENOR{} y a Moisés su siervo.
\end{verso}

\begin{filologia}{Éx 14:13--14, 31}
\heb{יְשׁוּעַת יְהוָה} & \heb{ישׁע} & ``la salvación de \SENOR'': obra de Dios, no de Israel. \\
\heb{יִלָּחֵם לָכֶם} & \heb{לחם} & ``peleará por vosotros'': Dios es el guerrero (\heb{אִישׁ מִלְחָמָה}, Éx 15:3). \\
\heb{תַּחֲרִישׁוּן} & \heb{חרשׁ} & ``callaréis, estaréis quietos'': la parte de Israel es no intervenir. \\
\heb{הַיָּד הַגְּדֹלָה} & \heb{יד} & literalmente ``la mano grande''; RVR1960: ``aquel grande hecho''. \\
\heb{וַיִּירְאוּ}\ldots{} \heb{וַיַּאֲמִינוּ} & \heb{ירא}, \heb{אמן} & ``temieron\ldots{} y creyeron'': ver el poder produce temor y fe. \\
\end{filologia}

El relato insiste en que Israel no hizo nada. Moisés extiende la mano, pero quien retira el mar es \SENOR. Y el fin del milagro no es el espectáculo, sino que el pueblo \emph{crea}: \heb{וַיַּאֲמִינוּ}, de la raíz de \emph{emuná}.

\subsection{El cántico del mar}

En la otra orilla, Israel canta. Es el primer canto de la Biblia, y lo dedica al poder de Dios:

\hebverso{יְמִינְךָ יְהוָה נֶאְדָּרִי בַּכֹּחַ יְמִינְךָ יְהוָה תִּרְעַץ אוֹיֵב\\ מִי־כָמֹכָה בָּאֵלִם יְהוָה מִי כָּמֹכָה נֶאְדָּר בַּקֹּדֶשׁ נוֹרָא תְהִלֹּת עֹשֵׂה פֶלֶא}
\begin{verso}{Éxodo 15:3, 6, 11, 18}
\SENOR{} es varón de guerra; \SENOR{} es su nombre\ldots{} Tu diestra, oh \SENOR, ha sido magnificada en poder; tu diestra, oh \SENOR, ha quebrantado al enemigo\ldots{} ¿Quién como tú, oh \SENOR, entre los dioses? ¿Quién como tú, magnífico en santidad, terrible en maravillosas hazañas, hacedor de prodigios?\ldots{} \SENOR{} reinará eternamente y para siempre.
\end{verso}

\begin{filologia}{Éx 15:6, 11, 18}
\heb{נֶאְדָּרִי בַּכֹּחַ} & \heb{אדר} & ``glorioso en poder''. El mismo verbo describe a Dios ``magnífico en santidad'' (\heb{נֶאְדָּר בַּקֹּדֶשׁ}, v.~11): poder y santidad son una sola gloria. \\
\heb{עֹשֵׂה פֶלֶא} & \heb{פלא} & ``hacedor de prodigio'': la raíz de Gn 18:14. Lo que en la tienda fue promesa, en el mar es hecho. \\
\heb{יִמְלֹךְ} & \heb{מלך} & ``reinará'' (v.~18): el poder de Dios no es un estallido, es un gobierno. \\
\end{filologia}

El cántico enlaza tres atributos: el poder (v.~6) se declara inseparable de la santidad (v.~11) y desemboca en el reinado (v.~18). La pregunta ``¿Quién como tú?'' es la manera bíblica de decir que el atributo es incomunicable, pues entre los \heb{אֵלִם}, los ``dioses'' o poderosos, no hay nadie igual.

\precision{en la liturgia judía el Cántico del mar se recita cada mañana, y ``¿Quién como tú?'' (Éx 15:11) y ``\SENOR{} reinará'' (Éx 15:18) se repiten en la bendición que sigue al \emph{Shemá}, la de la redención (Elbogen, 1993). Israel no recuerda el mar como un hecho del pasado, sino como el modelo de toda salvación.}

%===================================================================
\section{La omnipotencia en el judaísmo}\label{sec:judaismo-omni}

\subsection{“El que dijo al mundo: ¡basta!”}

Los sabios explicaron el nombre \emph{Shaddai} como \emph{she-dai}, ``el que basta'', y el Talmud pone esa lectura en boca de Resh Laquish:

\begin{quote}\small
\heb{מַאי דִּכְתִיב ״אֲנִי אֵל שַׁדַּי״ אֲנִי הוּא שֶׁאָמַרְתִּי לָעוֹלָם דַּי}\\
¿Qué significa lo que está escrito: ``Yo soy el Dios Todopoderoso''? Yo soy el que dijo al mundo: ¡Basta! (\emph{Jaguigá} 12a)
\end{quote}

El \emph{Midrash Rabá} conserva la misma explicación en nombre de R. Yitsjak: ``Yo soy el que dijo a mi mundo y a mis cielos: ¡Basta!''; sin esa palabra, la creación habría seguido extendiéndose (\emph{Bereshit Rabá} 46:3). Rashi, al comentar Génesis 17:1, da a la idea otra forma: ``Yo soy aquel cuya divinidad basta para toda criatura''. El Targum Ónkelos, que suele traducir los nombres y las imágenes de Dios, deja este intacto: \heb{אֲנָא אֵל שַׁדַּי}.

La explicación rabínica ilumina un aspecto del poder que la Biblia también conoce: Dios no solo crea, sino que \emph{pone límite}. ``Hasta aquí llegarás, y no pasarás adelante, y ahí parará el orgullo de tus olas'' (Job 38:11). El poder que dice ``sea'' es el mismo que dice ``basta''.

\precision{interpretación rabínica (midrash) del Nombre, no etimología. Coincide con la Septuaginta, que en Rut y en tres pasajes de Job traduce \emph{Shaddai} por \grc{ὁ ἱκανός}, ``el Suficiente''.}

\subsection{\emph{HaGevurá}: el Poder como nombre de Dios}

Los sabios emplean \heb{הַגְּבוּרָה} (\emph{haGevurá}, ``el Poder'') como forma reverente de nombrar a Dios. Al contar los 613 mandamientos, el Talmud afirma que Israel oyó dos de ellos directamente en el Sinaí:

\begin{quote}\small
\heb{״אָנֹכִי״ וְ״לֹא יִהְיֶה לְךָ״ מִפִּי הַגְּבוּרָה שְׁמַעְנוּם}\\
``Yo soy'' y ``No tendrás'' los oímos de la boca del Poder. (\emph{Makot} 23b--24a)
\end{quote}

Jesús recurre a ese mismo uso ante el sumo sacerdote:

\begin{verso}{Marcos 14:62}
Y Jesús le dijo: Yo soy; y veréis al Hijo del Hombre sentado a la diestra del poder de Dios, y viniendo en las nubes del cielo.
\end{verso}

El griego dice \grc{ἐκ δεξιῶν καθήμενον τῆς δυνάμεως}, ``sentado a la diestra del Poder''; la RVR1960 añade ``de Dios'' para aclarar el sentido. Es el uso que comentaba Charnock, para quien Dios y su poder son inseparables (Pink, 2020, cap. 9).

\precision{el paralelo es de lenguaje. El tratado \emph{Makot} es posterior al Nuevo Testamento, pero recoge una costumbre judía de llamar a Dios ``el Poder'' que Marcos ya atestigua en el siglo I.}

\subsection{La bendición \emph{Guevurot} de la Amidá}

La Amidá, la oración central del judaísmo, rezada tres veces al día, dedica su segunda bendición al poder de Dios. Se llama \emph{Guevurot}, ``poderes'':

\begin{quote}\small
\heb{אַתָּה גִּבּוֹר לְעוֹלָם אֲדֹנָי מְחַיֵּה מֵתִים אַתָּה רַב לְהוֹשִׁיעַ}\\
Tú eres poderoso para siempre, Señor; tú das vida a los muertos, grande eres para salvar\ldots{} Sostienes a los que caen, sanas a los enfermos y liberas a los cautivos\ldots{} Bendito eres tú, \SENOR, que das vida a los muertos. (Amidá, segunda bendición)
\end{quote}

La bendición define el poder de Dios por lo que hace con los débiles, y sus frases proceden de los salmos: ``Sostiene \SENOR{} a todos los que caen'' (Sal 145:14); ``\SENOR{} liberta a los cautivos; \SENOR{} abre los ojos a los ciegos; \SENOR{} levanta a los caídos'' (Sal 146:7--8). La cumbre de ese poder no es destruir, sino \emph{resucitar}: la bendición empieza y termina con \heb{מְחַיֵּה מֵתִים}, ``que das vida a los muertos''.

\precision{texto litúrgico según el sidur asquenazí (Sefaria). La estructura de la Amidá es antigua y la tradición la atribuye a la época de los sabios; su redacción exacta varió con los ritos (Elbogen, 1993).}

\subsection{Las tres llaves}

En esa misma bendición, durante el invierno, se menciona la lluvia: \emph{guevurot gueshamim}, ``los poderes de la lluvia''. El Talmud explica el nombre porque la lluvia ``cae con poder'', y en ese contexto R. Yojanán enseña:

\begin{quote}\small
\heb{שְׁלֹשָׁה מַפְתְּחוֹת בְּיָדוֹ שֶׁל הַקָּדוֹשׁ בָּרוּךְ הוּא שֶׁלֹּא נִמְסְרוּ בְּיַד שָׁלִיחַ}\\
\heb{מַפְתֵּחַ שֶׁל גְּשָׁמִים\quad מַפְתֵּחַ שֶׁל חַיָּה\quad וּמַפְתֵּחַ שֶׁל תְּחִיַּית הַמֵּתִים}\\
Tres llaves están en la mano del Santo, bendito sea, que no fueron entregadas a ningún enviado: la llave de la lluvia, la llave del parto y la llave de la resurrección de los muertos. (\emph{Taanit} 2a--b)
\end{quote}

Cada llave se prueba con un versículo. La de la lluvia, con Deuteronomio 28:12 (``Te abrirá \SENOR{} su buen tesoro, el cielo, para enviar la lluvia''); la del parto, con Génesis 30:22 (``se acordó Dios de Raquel\ldots{} y le concedió hijos''); la de la resurrección, con Ezequiel 37:13 (``cuando abra vuestros sepulcros''). Las tres abren lo que solo Dios puede abrir.

La historia de Génesis 17--18 reúne dos de esas llaves. Sara recibe la llave del parto, y Abraham, ``ya como muerto'', según dirá Pablo (Ro 4:19), recibe una vida que anticipa la de la resurrección. \emph{El Shaddai} es el Dios que tiene las llaves.

\subsection{“Cuya fuerza y poder llenan el mundo”}

El judaísmo no se detiene en el asombro ante el poder de Dios, sino que lo convierte en bendición. La Misná establece:

\begin{quote}\small
\heb{וְעַל הָרְעָמִים וְעַל הָרוּחוֹת אוֹמֵר בָּרוּךְ שֶׁכֹּחוֹ וּגְבוּרָתוֹ מָלֵא עוֹלָם}\\
Y por los truenos y por los vientos se dice: Bendito el que su fuerza y su poder llenan el mundo. (Misná, \emph{Berajot} 9:2)
\end{quote}

En la fórmula aparecen juntos los dos términos centrales de la familia del poder, \heb{כֹּחַ} y \heb{גְּבוּרָה}. Quien oye el trueno oye el salmo 29: ``voz de \SENOR{} con potencia''. El poder de Dios no se discute; se bendice cuando se manifiesta.

\subsection{El poder que se contiene}

El Talmud recoge una discusión sobre los tres adjetivos de Deuteronomio 10:17, ``el Dios grande, poderoso y temible''. Moisés los pronunció todos. Jeremías, al ver a los extranjeros destruir el templo, omitió ``temible'' (Jer 32:18). Daniel, al ver a Israel en el exilio, omitió ``poderoso'' (Dn 9:4), y preguntó:

\begin{quote}\small
\heb{גּוֹיִם מִשְׁתַּעְבְּדִים בְּבָנָיו אַיֵּה גְּבוּרוֹתָיו}\\
Las naciones esclavizan a sus hijos: ¿dónde están sus proezas? (\emph{Yomá} 69b)
\end{quote}

Los hombres de la Gran Asamblea restituyeron la fórmula completa (Neh 9:32) con una respuesta:

\begin{quote}\small
\heb{זוֹ הִיא גְּבוּרָתוֹ שֶׁכּוֹבֵשׁ אֶת יִצְרוֹ}\\
Esta es la fuerza de su fuerza: que domina su inclinación y es paciente con los malvados. (\emph{Yomá} 69b)
\end{quote}

La pregunta de Daniel es la de todo creyente que sufre: si Dios es todopoderoso, ¿por qué no actúa? Los sabios responden que el poder de Dios se muestra también en lo que \emph{no} hace todavía, porque se contiene por paciencia. El mismo R. Yojanán que habló de las tres llaves formuló otra regla:

\begin{quote}\small
\heb{כׇּל מָקוֹם שֶׁאַתָּה מוֹצֵא גְּבוּרָתוֹ שֶׁל הַקָּדוֹשׁ בָּרוּךְ הוּא אַתָּה מוֹצֵא עַנְוְתָנוּתוֹ}\\
Dondequiera que encuentras el poder del Santo, bendito sea, allí encuentras su humildad. (\emph{Meguilá} 31a)
\end{quote}

La prueba se toma de las tres partes de la Biblia hebrea. En la Torá, el ``Dios grande, poderoso y temible'' es el que ``hace justicia al huérfano y a la viuda'' (Dt 10:17--18); en los Profetas, el ``Alto y Sublime'' habita ``con el quebrantado y humilde de espíritu'' (Is 57:15); y en los Escritos, el que ``cabalga sobre los cielos'' es ``padre de huérfanos y defensor de viudas'' (Sal 68:4--5 [68:5--6]).

\precision{lecturas rabínicas. La idea de que el poder de Dios se manifiesta en paciencia y en cercanía a los humildes tiene apoyo bíblico (Nah 1:3: ``\SENOR{} es tardo para la ira y grande en poder''; Is 57:15), pero la formulación es de los sabios. Pablo dirá algo semejante con otra lógica: el poder de Dios se perfecciona en la debilidad (2~Co 12:9).}

\subsection{¿Quién es poderoso?}

Si la \emph{gevurá} de Dios consiste en contenerse, también la del ser humano. La Misná lo expresa en una sentencia célebre:

\begin{quote}\small
\heb{אֵיזֶהוּ גִבּוֹר הַכּוֹבֵשׁ אֶת יִצְרוֹ}\\
¿Quién es poderoso? El que domina su inclinación, como está dicho: ``Mejor es el que tarda en airarse que el fuerte; y el que se enseñorea de su espíritu, que el que toma una ciudad''. (\emph{Pirké Avot} 4:1; Pr 16:32)
\end{quote}

La omnipotencia, que es incomunicable, deja así un reflejo en el ser humano: el dominio de sí. Pablo lo llamará fruto del Espíritu, \grc{ἐγκράτεια}, ``templanza'' (Gá 5:23), palabra de la misma familia que \grc{κράτος}.

\subsection{Qumrán: el poder que resucita}

Entre los manuscritos de la cueva 4 de Qumrán figura un texto que los editores llamaron \emph{Apocalipsis mesiánico} (4Q521). Su fragmento 2, columna ii, describe lo que hará el Señor cuando ``los cielos y la tierra escuchen a su Mesías'': libertará a los prisioneros, devolverá la vista a los ciegos, enderezará a los encorvados, curará a los heridos, ``a los muertos hará vivir'' y anunciará buenas noticias a los humildes (traducción según García Martínez, 1992).

El texto sigue de cerca el salmo 146:7--8 y le añade la resurrección, el mismo recorrido de la bendición \emph{Guevurot}. Y cuando Juan el Bautista pregunta desde la cárcel si Jesús es el que había de venir, recibe una lista casi idéntica:

\begin{verso}{Lucas 7:22}
Y respondiendo Jesús, les dijo: Id, haced saber a Juan lo que habéis visto y oído: los ciegos ven, los cojos andan, los leprosos son limpiados, los sordos oyen, los muertos son resucitados, y a los pobres es anunciado el evangelio.
\end{verso}

\precision{edición crítica en Puech (1998, DJD XXV). El fragmento está dañado y su reconstrucción tiene partes discutidas; conviene cotejar la imagen en la Leon Levy Dead Sea Scrolls Digital Library. El paralelo con Lc 7:22 lo señalan muchos comentaristas, pero no prueba dependencia directa: muestra que en el judaísmo del siglo I la resurrección era la señal máxima del poder de Dios en los días del Mesías.}

%===================================================================
\section{La omnipotencia en Cristo, leída con Pablo}\label{sec:cristo-omni}

\subsection{“Su eterno poder”}

Pablo abre la carta a los Romanos con el poder de Dios, y lo hace en dos sentidos:

\begin{verso}{Romanos 1:16}
Porque no me avergüenzo del evangelio, porque es poder de Dios para salvación a todo aquel que cree; al judío primeramente, y también al griego.
\end{verso}

\begin{verso}{Romanos 1:20}
Porque las cosas invisibles de él, su eterno poder y deidad, se hacen claramente visibles desde la creación del mundo, siendo entendidas por medio de las cosas hechas, de modo que no tienen excusa.
\end{verso}

\begin{filologia}{Ro 1:16, 20}
\grc{δύναμις θεοῦ} & --- & ``poder de Dios'': el evangelio no solo habla del poder, \emph{es} poder que salva. \\
\grc{ἥ τε ἀΐδιος αὐτοῦ δύναμις} & --- & ``su poder eterno''. \grc{ἀΐδιος}, ``perpetuo''. Es lo que de Dios se ve en lo creado. \\
\grc{καθορᾶται} & --- & ``se ve claramente'': lo invisible se hace visible en lo hecho (\grc{τοῖς ποιήμασιν}). \\
\end{filologia}

Pablo lee la creación como Israel la había leído en los salmos y en Job: ``Los cielos cuentan la gloria de Dios'' (Sal 19:1). El poder que hizo el mundo es visible, y por eso nadie tiene excusa. Pero el mismo capítulo indica dónde se ve el poder que \emph{salva}: en el evangelio. El poder que creó es el que ahora recrea.

\subsection{Abraham: “poderoso para hacer”}

En Romanos 4, Pablo regresa al texto ancla de este capítulo y lee Génesis 17--18 como el retrato de la fe:

\begin{verso}{Romanos 4:17–21}
(como está escrito: Te he puesto por padre de muchas gentes) delante de Dios, a quien creyó, el cual da vida a los muertos, y llama las cosas que no son, como si fuesen. Él creyó en esperanza contra esperanza, para llegar a ser padre de muchas gentes, conforme a lo que se le había dicho: Así será tu descendencia. Y no se debilitó en la fe al considerar su cuerpo, que estaba ya como muerto (siendo de casi cien años), o la esterilidad de la matriz de Sara. Tampoco dudó, por incredulidad, de la promesa de Dios, sino que se fortaleció en fe, dando gloria a Dios, plenamente convencido de que era también poderoso para hacer todo lo que había prometido.
\end{verso}

\begin{center}\small\renewcommand{\arraystretch}{1.25}
\begin{tabular}{@{}L{4.4cm}L{5.0cm}L{5.0cm}@{}}
\toprule
\textbf{Génesis} & \textbf{Romanos 4} & \textbf{Comentario} \\ \midrule
Abram, 99 años (17:1) & ``siendo de casi cien años'' (v.~19) & Pablo conserva el dato \\
Sara estéril (18:11) & ``la esterilidad de la matriz de Sara'' & la llave del parto \\
\emph{El Shaddai} (17:1) & ``el cual da vida a los muertos'' (v.~17) & el Nombre, leído como poder de resurrección \\
\heb{הֲיִפָּלֵא מֵיְהוָה דָּבָר} (18:14) & \grc{δυνατός ἐστιν καὶ ποιῆσαι} (v.~21) & la pregunta, convertida en certeza \\
\heb{תָמִים}, ``sé perfecto'' (17:1) & ``se fortaleció en fe, dando gloria a Dios'' (v.~20) & la respuesta humana \\
\bottomrule
\end{tabular}\end{center}

La fórmula ``el cual da vida a los muertos'' (\grc{τοῦ ζῳοποιοῦντος τοὺς νεκρούς}) resuena con la bendición judía \heb{מְחַיֵּה מֵתִים}. Pablo lee el nacimiento de Isaac como una resurrección anticipada, puesto que nace de dos cuerpos ``como muertos''. En los versículos siguientes extrae la consecuencia: esa misma fe es la de quienes creen ``en el que levantó de los muertos a Jesús, Señor nuestro'' (Ro 4:24). La fe de Abraham es fe en la omnipotencia.

\precision{el parecido entre Ro 4:17 y la bendición \emph{Guevurot} es de lenguaje y de fe compartida en la resurrección; no se puede afirmar que Pablo cite la Amidá, cuya redacción exacta es posterior. ``Llama las cosas que no son, como si fuesen'' recoge la fe judía en la creación por la palabra (Gn 1; Sal 33:9).}

\subsection{Cristo, poder de Dios}

En Corinto, Pablo se enfrenta a dos expectativas sobre el poder: los judíos piden señales y los griegos buscan sabiduría (1~Co 1:22). A ambas responde con un crucificado:

\begin{verso}{1 Corintios 1:18}
Porque la palabra de la cruz es locura a los que se pierden; pero a los que se salvan, esto es, a nosotros, es poder de Dios.
\end{verso}

\begin{verso}{1 Corintios 1:24–25}
mas para los llamados, así judíos como griegos, Cristo poder de Dios, y sabiduría de Dios. Porque lo insensato de Dios es más sabio que los hombres, y lo débil de Dios es más fuerte que los hombres.
\end{verso}

\grc{Χριστὸν θεοῦ δύναμιν}: Cristo no solo tiene poder, sino que \emph{es} el poder de Dios. Quien esperaba un milagro como el del mar Rojo encuentra el poder de Dios en la cruz, justo donde parecía faltar. Pablo cierra el argumento con el verso que une esta obra: ``El que se gloría, gloríese en el Señor'' (1~Co 1:31; Jer 9:24). Ante la cruz se quedan sin gloria el sabio y el \emph{valiente}.

\subsection{La supereminente grandeza de su poder}

Si la cruz es el poder escondido, la resurrección es el poder manifestado. En Efesios, Pablo ora para que los creyentes conozcan

\begin{verso}{Efesios 1:19–21}
y cuál la supereminente grandeza de su poder para con nosotros los que creemos, según la operación del poder de su fuerza, la cual operó en Cristo, resucitándole de los muertos y sentándole a su diestra en los lugares celestiales, sobre todo principado y autoridad y poder y señorío, y sobre todo nombre que se nombra, no sólo en este siglo, sino también en el venidero.
\end{verso}

\begin{filologia}{Ef 1:19--20}
\grc{ὑπερβάλλον μέγεθος} & --- & ``grandeza que sobrepasa'': Pablo agota el idioma. \\
\grc{δύναμις} & --- & ``poder, capacidad''. \\
\grc{ἐνέργεια} & --- & ``operación, eficacia'': el poder en acto. \\
\grc{κράτος} & --- & ``dominio''. \\
\grc{ἰσχύς} & --- & ``fuerza''. \\
\end{filologia}

Pablo acumula cuatro palabras de poder para una sola obra: resucitar a Cristo y sentarlo a la diestra, según el salmo 110:1. La resurrección de Cristo es para él la medida de la omnipotencia, y esa misma fuerza obra ``para con nosotros los que creemos''. De ahí que la carta pueda cerrar su primera parte con esta doxología:

\begin{verso}{Efesios 3:20–21}
Y a Aquel que es poderoso para hacer todas las cosas mucho más abundantemente de lo que pedimos o entendemos, según el poder que actúa en nosotros, a él sea gloria en la iglesia en Cristo Jesús por todas las edades, por los siglos de los siglos. Amén.
\end{verso}

\precision{la autoría paulina de Efesios se discute. La misma relación entre el poder de Dios y la resurrección se encuentra en cartas indiscutidas: ``Dios, que levantó al Señor, también a nosotros nos levantará con su poder'' (1~Co 6:14); ``se siembra en debilidad, resucitará en poder'' (1~Co 15:43).}

\subsection{Mi poder se perfecciona en la debilidad}

Pablo pidió tres veces que le fuera quitado el ``aguijón en la carne'', y la respuesta de Cristo transformó su manera de entender el poder:

\begin{verso}{2 Corintios 12:9–10}
Y me ha dicho: Bástate mi gracia; porque mi poder se perfecciona en la debilidad. Por tanto, de buena gana me gloriaré más bien en mis debilidades, para que repose sobre mí el poder de Cristo. Por lo cual, por amor a Cristo me gozo en las debilidades, en afrentas, en necesidades, en persecuciones, en angustias; porque cuando soy débil, entonces soy fuerte.
\end{verso}

\begin{filologia}{2~Co 12:9}
\grc{ἀρκεῖ σοι ἡ χάρις μου} & --- & ``te basta mi gracia''. \grc{ἀρκέω}, ``bastar, ser suficiente''. \\
\grc{ἐν ἀσθενείᾳ τελεῖται} & --- & ``en debilidad se lleva a término'': la debilidad es el lugar donde el poder llega a su fin. \\
\grc{ἐπισκηνώσῃ} & --- & ``haga tienda sobre'' (de \grc{σκηνή}, ``tienda''): el poder de Cristo acampa sobre el débil, como la gloria sobre el tabernáculo. \\
\end{filologia}

La palabra de Cristo a Pablo dice lo mismo que la tradición judía decía de \emph{Shaddai}: es suficiente. El Dios ``que basta'' le asegura a Pablo que su gracia \emph{basta}. Y el poder que se perfecciona en la debilidad es el de un Cristo que ``fue crucificado en debilidad'', pero ``vive por el poder de Dios'' (2~Co 13:4).

\precision{relacionar \grc{ἀρκεῖ} con la lectura rabínica \emph{she-dai} es una lectura propia de esta serie, no una alusión explícita de Pablo. Es un eco de sentido, no de palabras.}

\subsection{“Dice el Señor Todopoderoso”}

Pablo usa \grc{παντοκράτωρ} una sola vez, dentro de una cadena de citas del Antiguo Testamento:

\begin{verso}{2 Corintios 6:18}
Y seré para vosotros por Padre, y vosotros me seréis hijos e hijas, dice el Señor Todopoderoso.
\end{verso}

La cita procede de la promesa de Natán a David: ``Así ha dicho \SENOR{} de los ejércitos\ldots{} Yo le seré a él padre, y él me será a mí hijo'' (2~S 7:8, 14), donde la Septuaginta traduce \heb{יְהוָה צְבָאוֹת} como \grc{κύριος παντοκράτωρ}. Pablo extiende a la iglesia, ``hijos e hijas'', la promesa hecha al hijo de David. El Todopoderoso no se presenta aquí como amenaza, sino como \emph{Padre}.

\subsection{Nada frustra su voluntad}

La segunda mitad de la definición ---nada frustra la voluntad de Dios--- es para Pablo fuente de seguridad. Anticipa la objeción, ``¿quién ha resistido a su voluntad?'' (Ro 9:19), y la responde con la confesión de Isaías y Daniel. Su conclusión pastoral, sin embargo, está en Romanos 8:

\begin{verso}{Romanos 8:38–39}
Por lo cual estoy seguro de que ni la muerte, ni la vida, ni ángeles, ni principados, ni potestades, ni lo presente, ni lo por venir, ni lo alto, ni lo profundo, ni ninguna otra cosa creada nos podrá separar del amor de Dios, que es en Cristo Jesús Señor nuestro.
\end{verso}

Entre los poderes incapaces de separar del amor de Dios, Pablo nombra las \grc{δυνάμεις}, las ``potestades''. Ningún poder creado puede vencer al poder del Creador, y ese poder está al servicio del amor. Por eso la esperanza final es la resurrección del cuerpo:

\begin{verso}{Filipenses 3:21}
el cual transformará el cuerpo de la humillación nuestra, para que sea semejante al cuerpo de la gloria suya, por el poder con el cual puede también sujetar a sí mismo todas las cosas.
\end{verso}

El último enemigo, la muerte, ``será destruido'' (1~Co 15:26). La llave de la resurrección, que según los sabios no fue entregada a ningún enviado, está en las manos de Cristo.

\subsection{La plenitud del poder en Cristo}

Pablo reúne en Cristo todos los escenarios del poder. En él fueron creadas todas las cosas (Col 1:16) y en él subsisten (Col 1:17); en él se mostró la supereminente grandeza del poder al resucitar (Ef 1:19--20), y por él todas las cosas serán sujetadas (Flp 3:21; 1~Co 15:24--28). El niño de Isaías se llamaría \heb{אֵל גִּבּוֹר}, ``Dios Fuerte'' (Is 9:6 [9:5]); el Crucificado es ``poder de Dios'' (1~Co 1:24). ``En él habita corporalmente toda la plenitud de la Deidad'' (Col 2:9): el poder que hizo nacer a Isaac, abrió el mar y resucita a los muertos tiene rostro.

\precision{leer \emph{El Gibbor} de Is 9:6 [9:5] como título de Cristo es una lectura mesiánica cristiana; en su contexto es el nombre del rey davídico prometido, y la Septuaginta lo traduce de otro modo (``ángel del gran consejo''). Colosenses es de autoría discutida; la mediación de Cristo en la creación está también en 1~Co 8:6.}

%===================================================================
\section{Vivir ante el Todopoderoso}

La omnipotencia de Dios no es solo una doctrina que se afirma; reclama respuestas concretas.

\begin{enumerate}
\item \textbf{Adorar.} Ante el poder de Dios, la primera reacción de Israel fue cantar: ``¿Quién como tú, oh \SENOR, entre los dioses?'' (Éx 15:11). La última escena de la Biblia repite ese gesto: ``Santo, santo, santo es el Señor Dios Todopoderoso'' (Ap 4:8). El poder de Dios solo se conoce de verdad cuando se adora.

\item \textbf{Confiar en lo imposible.} La pregunta de Génesis 18:14 sigue abierta para cada creyente. Hay promesas de Dios que hemos dado por imposibles, como Sara detrás de la puerta, y la omnipotencia invita a orar sin achicar a Dios: ``a Aquel que es poderoso para hacer todas las cosas mucho más abundantemente de lo que pedimos o entendemos'' (Ef 3:20).

\item \textbf{Gloriarse en el Señor.} ``Ni en su valentía se alabe el valiente'' (Jer 9:23). La fuerza, la salud, la influencia y la capacidad de trabajo son dones recibidos, porque ``él te da el poder'' (Dt 8:18). Gloriarse en ellos es olvidar de dónde vienen.

\item \textbf{Recibir fuerza en la debilidad.} El poder de Dios no solo está por encima del ser humano; también se entrega al que no lo tiene:

\begin{verso}{Isaías 40:29–31}
El da esfuerzo al cansado, y multiplica las fuerzas al que no tiene ningunas. Los muchachos se fatigan y se cansan, los jóvenes flaquean y caen; pero los que esperan a \SENOR{} tendrán nuevas fuerzas; levantarán alas como las águilas; correrán, y no se cansarán; caminarán, y no se fatigarán.
\end{verso}

Así lo vivió Pablo: ``cuando soy débil, entonces soy fuerte'' (2~Co 12:10).

\item \textbf{Dominarse.} El reflejo humano de la \emph{gevurá} consiste en contenerse como Dios se contiene. Lo enseñan los Proverbios, ``Mejor es el que tarda en airarse que el fuerte'' (Pr 16:32); los sabios, ``¿quién es poderoso? El que domina su inclinación'' (\emph{Avot} 4:1); y Pablo, que llama a esto ``templanza'' (Gá 5:23). El creyente no imita la omnipotencia de Dios ejerciendo poder sobre otros, sino sobre sí mismo.
\end{enumerate}

Queda una última palabra para quien sufre. Noemí dijo: ``en grande amargura me ha puesto el Todopoderoso'' (Rt 1:20), y Job habló de las ``saetas del Todopoderoso'' (Job 6:4). La Biblia permite al creyente hablar así de \emph{Shaddai}. Pero el libro de Rut termina con un niño en brazos de Noemí, del que saldrá David (Rt 4:16--17), y el de Job con una visión de Dios (Job 42:5). El mismo poder que parecía aplastar estaba escribiendo una historia mayor.

%===================================================================
\section{Síntesis}

\begin{center}
\begin{tcolorbox}[colback=papiro,colframe=azul,boxrule=1pt,arc=3pt,width=0.95\linewidth]
Dios es \emph{El Shaddai}: el que basta, el que hace nacer lo que no puede nacer (Gn 17:1; 18:14), el que abre el mar (Éx 14--15) y dice ``sea'', y es (Gn 1:3; Sal 33:9). Todo lo que quiere, lo hace (Sal 115:3), y nada frustra su consejo (Job 42:2; Is 46:10). Su poder es uno con su misericordia (Sal 62:11--12) y con su santidad (Éx 15:11). Israel lo bendice en el trueno y en la Amidá, y los sabios lo reconocen en su paciencia y en su humildad. Pablo lo contempla en Cristo crucificado y resucitado, ``poder de Dios'' (1~Co 1:24; Ef 1:19--20), y lo recibe como gracia que basta en la debilidad (2~Co 12:9).
\end{tcolorbox}
\end{center}

\begin{center}\small\renewcommand{\arraystretch}{1.25}
\begin{tabular}{@{}L{2.6cm}L{11.8cm}@{}}
\toprule
\textbf{Mirada} & \textbf{Aporte al atributo} \\ \midrule
Torá & \emph{El Shaddai} y la descendencia; \heb{הֲיִפָּלֵא} (Gn 18:14); la palabra creadora (Gn 1); el mar Rojo (Éx 14--15); ``grande, poderoso y temible'' (Dt 10:17); ``yo hago morir y hago vivir'' (Dt 32:39). \\
Salmos & ``de Dios es el poder'' y la misericordia (62:11--12); ``él dijo, y fue hecho'' (33:9); ``todo lo que quiso ha hecho'' (115:3); la voz, el brazo, el guerrero y la sombra (29; 89; 24; 91). \\
Judaísmo & \emph{she-dai} (\emph{Jaguigá} 12a; \emph{Bereshit Rabá} 46:3; Rashi); \emph{haGevurá} (\emph{Makot} 24a); la Amidá; las tres llaves (\emph{Taanit} 2a--b); \emph{Berajot} 9:2; \emph{Yomá} 69b; \emph{Meguilá} 31a; \emph{Avot} 4:1; 4Q521. \\
Teológica & el poder da vida a los atributos (Charnock, en Pink, 2020, cap. 9); el poder escondido (Hab 3:4; Job 26:14); creación, preservación, gobierno y juicio; temblar, adorar y confiar. \\
Paulina & Ro 1:16, 20; 4:17--21; 8:38--39; 1~Co 1:18--25; 2~Co 6:18; 12:9--10; 13:4; Ef 1:19--21; 3:20--21; Flp 3:21. \\
\bottomrule
\end{tabular}\end{center}

\begin{verso}{Génesis 18:14}
¿Hay para Dios alguna cosa difícil? Al tiempo señalado volveré a ti, y según el tiempo de la vida, Sara tendrá un hijo.
\end{verso}

El verso que une esta obra quita la gloria al valiente para dársela al que conoce a Dios. Y conocer a Dios es conocer al único poderoso, el que usa su poder para dar vida.

%===================================================================
\section*{Glosario del capítulo}

\begin{description}\setlength{\itemsep}{2pt}
\item[\emph{Adynatéō}] (\grc{ἀδυνατέω}) ser imposible (Gn 18:14 LXX; Lc 1:37).
\item[\emph{Batsar}] (\heb{בָּצַר}) en nifal, ser impedido, resultar inalcanzable (Gn 11:6; Job 42:2).
\item[\emph{Dýnamis}] (\grc{δύναμις}) poder, capacidad (Ro 1:16, 20; 1~Co 1:24).
\item[\emph{El Shaddai}] (\heb{אֵל שַׁדַּי}) Dios Todopoderoso (Gn 17:1; Éx 6:3).
\item[\emph{Enérgeia}] (\grc{ἐνέργεια}) operación eficaz, poder en acto (Ef 1:19; Flp 3:21).
\item[\emph{Gevurá}] (\heb{גְּבוּרָה}) fuerza, proeza; en los sabios, nombre de Dios (\emph{haGevurá}).
\item[\emph{Gibbor}] (\heb{גִּבּוֹר}) poderoso, valiente (Dt 10:17; Is 9:6 [9:5]; Jer 9:23 [9:22]).
\item[\emph{Hikanós}] (\grc{ἱκανός}) suficiente; versión de \emph{Shaddai} en la LXX de Rut y Job.
\item[\emph{Ischýs}] (\grc{ἰσχύς}) fuerza (Ef 1:19).
\item[\emph{Kóaj}] (\heb{כֹּחַ}) vigor, capacidad (Éx 15:6; Jer 32:17).
\item[\emph{Krátos}] (\grc{κράτος}) dominio (Ef 1:19).
\item[\emph{Oz}] (\heb{עֹז}) fuerza firme, fortaleza (Sal 62:11 [62:12]).
\item[\emph{Palá}] (\heb{פָּלָא}) ser extraordinario, demasiado difícil (Gn 18:14; Jer 32:17).
\item[\emph{Pantokrátōr}] (\grc{παντοκράτωρ}) el que lo domina todo; Todopoderoso (2~Co 6:18; Ap 1:8).
\item[\emph{Yajol}] (\heb{יָכֹל}) poder, ser capaz (Job 42:2).
\item[\emph{Zeróa netuyá}] (\heb{זְרוֹעַ נְטוּיָה}) brazo extendido (Éx 6:6; Dt 26:8).
\end{description}

%===================================================================
